% Схема алгоритма декомпозиции дальней цели на этапы кампании.
\begin{tikzpicture}[
  font=\footnotesize,
  node distance=3.6mm,
  term/.style={draw, rounded rectangle, align=center, minimum width=3.4cm, minimum height=8mm},
  proc/.style={draw, align=center, minimum width=7.0cm, minimum height=9mm, inner sep=3pt},
  sub/.style={proc, inner xsep=4mm, append after command={
    \pgfextra{
      \draw ([xshift=2mm]\tikzlastnode.north west) -- ([xshift=2mm]\tikzlastnode.south west);
      \draw ([xshift=-2mm]\tikzlastnode.north east) -- ([xshift=-2mm]\tikzlastnode.south east);
    }}},
  cond/.style={draw, diamond, aspect=3.4, align=center, inner sep=0pt, minimum width=7.0cm},
  exit/.style={draw, rounded rectangle, align=center, text width=3.2cm, minimum height=9mm, inner sep=3pt},
  side/.style={draw, align=center, text width=3.4cm, minimum height=9mm, inner sep=3pt},
  arr/.style={-{Stealth[length=1.8mm]}, semithick},
  ylbl/.style={font=\scriptsize, above, pos=0.3},
  nlbl/.style={font=\scriptsize, right, pos=0.4},
]

\node[term] (start) {Начало: цель, инвентарь, режим};
\node[sub, below=of start] (plan) {Построить план цели от инвентаря\\(рисунок~\ref{fig:planner})};
\node[cond, below=of plan] (feasible) {План выполним?};
\node[proc, below=of feasible] (marks) {Обойти дерево плана снизу вверх и выделить вехи:\\инструменты добычи, оборудование рецептов, цель};
\node[proc, below=of marks] (next) {Взять очередную веху; состояние~--~прогноз\\инвентаря после принятых этапов};
\node[sub, below=of next] (check) {Проверить веху как задание: план от прогноза,\\трудоёмкость по формуле~eqref{eq:quest}};
\node[cond, below=of check] (fits) {Трудоёмкость\\в допустимых пределах?};
\node[proc, below=of fits] (accept) {Принять этап; прогноз: добавить цель этапа,\\вычесть израсходованное его планом};
\node[cond, below=of accept] (more) {Вехи остались?};
\node[proc, below=of more] (rewards) {Назначить награды этапов по правилу отбора};
\node[term, below=of rewards] (end) {Конец: цепочка этапов с наградами};

\node[exit, right=10mm of feasible] (unreach) {Отказ: цель\\недостижима};
\node[cond, right=8mm of fits, minimum width=3.9cm, aspect=2.4] (split) {Этап делится?};
\node[side, above=7mm of split] (heavy) {Выделить самое тяжёлое поддерево материала в предшествующий этап};
\node[exit, below=7mm of split] (fail) {Отказ: этап\\и причина};

\foreach \a/\b in {start/plan, plan/feasible, marks/next, next/check, check/fits, accept/more, rewards/end}
  \draw[arr] (\a) -- (\b);
\draw[arr] (feasible) -- node[nlbl] {да} (marks);
\draw[arr] (feasible) -- node[ylbl] {нет} (unreach);
\draw[arr] (fits) -- node[nlbl] {да} (accept);
\draw[arr] (fits) -- node[ylbl] {нет} (split);
\draw[arr] (split) -- node[nlbl] {да} (heavy);
\draw[arr] (split) -- node[nlbl] {нет} (fail);
\draw[arr] (heavy.north) |- (check.east);
\draw[arr] (more) -- node[nlbl] {нет} (rewards);
% Возврат к следующей вехе по линии слева.
\draw[arr] (more.west) -- node[font=\scriptsize, above, pos=0.5] {да} ++(-9mm, 0) |- (next.west);

\end{tikzpicture}
